\documentclass[11pt]{article}
\usepackage[margin=1in]{geometry}
\usepackage{amsmath,amssymb,amsthm}
\usepackage{booktabs}
\usepackage{hyperref}
\hypersetup{colorlinks=true,linkcolor=blue,urlcolor=blue}

\newcommand{\eps}{\varepsilon}
\newcommand{\epsth}{\eps_\theta}
\newcommand{\epsfr}{\eps_{\bar\theta}}
\newcommand{\xs}{x_\sigma}
\newcommand{\E}{\mathbb{E}}
\newcommand{\R}{\mathbb{R}}
\newcommand{\norm}[1]{\left\lVert #1 \right\rVert}
\newcommand{\cnull}{\varnothing}
\newcommand{\cpl}{c_{\mathrm{pl}}}

\theoremstyle{plain}
\newtheorem{proposition}{Proposition}
\newtheorem{corollary}[proposition]{Corollary}
\theoremstyle{remark}
\newtheorem{remark}[proposition]{Remark}

\title{Four correction losses for product-of-experts composition,\\ written in one notation}
\author{poe\_repair\_min, thread off \texttt{designing-the-correction-loss} claim 4}
\date{2026-09-12}

\begin{document}
\maketitle

\begin{abstract}
Four proposed training objectives for the SDXL product-of-experts corrector are written in a
single notation so they can be compared as mathematics rather than as descriptions. The
comparison reduces to one fact: all four sit inside the same additive family
$T = u + f_i + f_j$, so at a state that does not identify the pair they share one least-squares
floor, and two of the four are strict subfamilies with a strictly worse floor. The document
also settles whether the plurality-conditioned composition is the score of a density (it is),
what that density is, whether prompt concatenation is well defined on SDXL's two text encoders
(only at the string level), and what the objective costs per step (less than the objective it
replaces).
\end{abstract}

\tableofcontents

\section{Setup and conventions}

\paragraph{The forward process.}
Write $x \sim p_{\mathrm{data}}$ for a clean latent, $z \sim \mathcal{N}(0,I)$, and
$\xs = x + \sigma z$ for the noised state at scale $\sigma$. The network
$\epsth(\xs,\sigma,c)$ predicts noise. Frozen (pre-adapter) weights are written $\bar\theta$,
so $\epsfr$ is the base SDXL prediction and $\epsth$ the adapted one.

\paragraph{Score and noise conventions, and one sign to check.}
The two spellings are related by
\begin{equation}
  \nabla_{\xs} \log p_\sigma(\xs) \;=\; -\,\frac{1}{\sigma}\,\E[\,z \mid \xs\,],
  \qquad\text{so}\qquad
  s_\theta(\xs,\sigma,c) \;:=\; \nabla \log p_\sigma(\xs \mid c) \;\approx\; -\frac{\epsth}{\sigma}.
  \label{eq:convention}
\end{equation}
The score-space and noise-space losses
\begin{equation*}
  \norm{\tfrac{z}{\sigma} - (s_1 + s_2 - s_0)}^2
  \qquad\text{and}\qquad
  \sigma^{-2}\norm{z - (\eps_1 + \eps_2 - \eps_0)}^2
\end{equation*}
agree \emph{only if} the symbol $s$ denotes $\eps/\sigma$, minus the score of
\eqref{eq:convention}. With
$s$ read as the score itself the signs invert. The two forms differ by $\sigma^{-2}$, which is
absorbed into the per-level weight $\lambda(\sigma)$ below, so nothing downstream depends on
the choice. Everything from here on is in noise space.

\paragraph{The composition rule and guidance.}
The composed prediction and the guided sampler step are
\begin{equation}
  T \;=\; \eps_1 + \eps_2 - \eps_\cnull,
  \qquad
  T^{(w)} \;=\; \eps_\cnull + w(\eps_1 - \eps_\cnull) + w(\eps_2 - \eps_\cnull),
  \label{eq:compose}
\end{equation}
implemented as \texttt{poe\_eps} and \texttt{guided\_eps} in
\texttt{poe\_repair/\_sdxl/metrics.py:12--17} and combined in
\texttt{poe\_repair/experiments/one\_pair\_one\_seed/trainer.py:331--333} at $w = 7.5$.
Setting $w = 1$ in $T^{(w)}$ returns $T$, so guidance is a dial on one expression.

\paragraph{Branch coordinates.}
Fix a state. Write
\begin{equation}
  u \;:=\; \epsth(\xs,\sigma,\cnull),
  \qquad
  f_i \;:=\; \epsth(\xs,\sigma,c_i) - u,
  \label{eq:coords}
\end{equation}
so that, for the pair $(i,j)$,
\begin{equation}
  T_{ij} \;=\; u + f_i + f_j,
  \qquad
  T^{(w)}_{ij} \;=\; u + w\,(f_i + f_j)
  \;=\; w(\eps_i + \eps_j) - (2w-1)\,u .
  \label{eq:Tij}
\end{equation}
The right-hand form of \eqref{eq:Tij} records that the null branch is read twice at inference,
with net coefficient $-(2w-1) = -14$ at $w = 7.5$, against $-1$ in the unguided $T$.

\paragraph{The per-level weight.}
Every variant carries an optional weight $\lambda(\sigma)$ averaged over $L$ levels. This maps
onto \texttt{\_loss\_weights} at \texttt{trainer.py:128--151}, which implements
$\lambda_{\eps} \equiv 1$ and $\lambda_{x_0} = \min\!\big((1-\bar\alpha_t)/\bar\alpha_t,\,
\mathrm{cap}\big)$ normalised to unit mean over the DDIM grid. It plays no role in the
comparison below, since it multiplies every variant identically.

\section{The objective being replaced, and the objective proposed}

\paragraph{Today: distillation onto the model's own joint-prompt branch.}
\begin{equation}
  \mathcal{L}_{\mathrm{distill}}
  \;=\;
  \E_{(\xs,\sigma)\sim \text{cache}}\,
  \lambda(\sigma)\,
  \norm{\,T^{(w)}_{ij}[\theta] \;-\; \epsfr^{(w)}(\xs,\sigma,c_{i\wedge j})\,}^2 ,
  \label{eq:distill}
\end{equation}
with $c_{i\wedge j}$ the literal joint string \texttt{"a cat and a dog"} produced by
\texttt{joint\_prompt} and the states $\xs$ drawn from cached PoE sampling trajectories.

\paragraph{Proposed: denoising score matching against co-occurrence data.}
\begin{equation}
  \boxed{\;
  \mathcal{L}_{\mathrm{DSM}}
  \;=\;
  \frac{1}{L}\sum_{\ell=1}^{L} \lambda(\sigma_\ell)\;
  \E_{\substack{x \sim p_{\mathrm{co}}\\ z \sim \mathcal{N}(0,I)}}
  \norm{\,z \;-\; \big[\,\epsth(\xs,\sigma_\ell,c_1) + \epsth(\xs,\sigma_\ell,c_2)
  - \epsth(\xs,\sigma_\ell,\cnull)\,\big]\,}^2 \;}
  \label{eq:dsm}
\end{equation}
where $p_{\mathrm{co}}$ is a distribution over real images in which both concepts appear as
separate objects.

\begin{proposition}[What \eqref{eq:dsm} fits]
\label{prop:minimiser}
For each $\sigma$, the unconstrained minimiser of $\E\norm{z - g(\xs)}^2$ over measurable $g$
is $g^\star(\xs) = \E[z \mid \xs] = -\sigma \nabla \log p_{\mathrm{co},\sigma}(\xs)$. Over a
constrained family $\mathcal{F}$ of achievable composed predictions, the minimiser is the
$L^2(p_{\mathrm{co},\sigma})$ projection of $g^\star$ onto $\mathcal{F}$.
\end{proposition}

\begin{remark}
This is the part of the proposal that is doing real work, and it is not the change of target
images. The composed field of \eqref{eq:compose} is the score of a product at $\sigma = 0$
only; the score of a noised product is not the sum of noised scores, which is Du et al.'s
objection to composed sampling (arXiv:2302.11552). Objective \eqref{eq:dsm} regresses at every
$\sigma_\ell$ separately, so it fits the composition \emph{at each noise level} instead of
assuming it transports from clean data.
\end{remark}

\section{The four variants}

All four share the target $z$, the weight $\lambda(\sigma)$, and the expectation of
\eqref{eq:dsm}. They differ in what is conditioned on and what is free. Write $\bar\theta$ for
a frozen branch.

\paragraph{(a) Plain: three adapted branches.}
\begin{equation}
  T^{(a)}_{ij} \;=\; \epsth(\xs,\sigma,c_i) + \epsth(\xs,\sigma,c_j) - \epsth(\xs,\sigma,\cnull)
  \;=\; u + f_i + f_j .
  \label{eq:va}
\end{equation}

\paragraph{(b) Plurality-conditioned base measure.}
With $\cpl$ a plurality phrase and $\oplus$ string concatenation before encoding
(Section~\ref{sec:concat}),
\begin{equation}
  T^{(b)}_{ij} \;=\; \epsth(\xs,\sigma,c_i \oplus \cpl) + \epsth(\xs,\sigma,c_j \oplus \cpl)
  - \epsth(\xs,\sigma,\cpl)
  \;=\; u_{\mathrm{pl}} + g_i + g_j ,
  \label{eq:vb}
\end{equation}
where $u_{\mathrm{pl}} := \epsth(\xs,\sigma,\cpl)$ and
$g_i := \epsth(\xs,\sigma,c_i \oplus \cpl) - u_{\mathrm{pl}}$.

\paragraph{(c) Null-only adapter.}
Concept branches and the target frozen, one adapted base:
\begin{equation}
  T^{(c)}_{ij} \;=\; \epsfr(\xs,\sigma,c_i) + \epsfr(\xs,\sigma,c_j) - \epsth(\xs,\sigma,\cpl)
  \;=\; \bar u + \bar f_i + \bar f_j \;-\; \delta ,
  \label{eq:vc}
\end{equation}
with $\delta := \epsth(\xs,\sigma,\cpl) - \epsfr(\xs,\sigma,\cnull)$.
The free object is the single vector field $\delta$, shared by every pair.

\paragraph{(d) Connective adapter.} Two inequivalent readings.
\begin{align}
  \text{(d-lora)}\quad
  T^{(d\ell)}_{ij} &\;=\; \epsth(\xs,\sigma,\text{``}c_i\ \mathrm{and}\text{''})
  + \epsth(\xs,\sigma,\text{``}c_j\ \mathrm{and}\text{''})
  - \epsth(\xs,\sigma,\text{``and''}),
  \label{eq:vdl}\\[2pt]
  \text{(d-tok)}\quad
  T^{(dt)}_{ij} &\;=\; \epsfr(\xs,\sigma,[c_i, v]) + \epsfr(\xs,\sigma,[c_j, v])
  - \epsfr(\xs,\sigma,[v]),
  \qquad v \in \R^{d_{\mathrm{text}}} \text{ free}.
  \label{eq:vdt}
\end{align}
In \eqref{eq:vdl} the free parameters are LoRA factors on the cross-attention projections,
which act on every token's key and value, so the adapter is not localised to the word
\emph{and}; only the branch prompts change. In \eqref{eq:vdt} the free parameter is a single
learned token embedding $v$ inserted into all three prompts, with the network frozen, which is
the Textual Inversion parameterisation (arXiv:2208.01618). These are different families and
must not share a row.

\section{The common skeleton, and why capacity does not separate (a), (b) and (d-lora)}

\paragraph{Every variant is additive in the two concepts.}
Each of \eqref{eq:va}--\eqref{eq:vdt} has the form
\begin{equation}
  T_{ij} \;=\; m \;+\; \phi_i \;+\; \phi_j
  \label{eq:additive}
\end{equation}
for some shared $m$ and some per-concept $\phi$. This is forced by the architecture rather than
chosen: each branch is conditioned on one concept, and the combination rule \eqref{eq:compose}
is a fixed linear map with no term reading both concepts. Equation~\eqref{eq:additive} is a
two-way model with no interaction.

\paragraph{Gauge freedom.}
\begin{proposition}[The additive parameterisation is not identified]
\label{prop:gauge}
Fix a state and let $G$ be the graph on concepts with an edge per trained pair. The
reparameterisation $m \mapsto m + c$, $\phi_i \mapsto \phi_i + a_i$ leaves every $T_{ij}$ in
\eqref{eq:additive} unchanged if and only if $a_i + a_j + c = 0$ for every edge $(i,j)$. On a
connected component that is not bipartite the only solutions are $a_i \equiv -c/2$, a
one-dimensional family per component; on a bipartite component there is one extra dimension
($a \equiv \alpha$ on one side, $\beta$ on the other, $\alpha + \beta + c = 0$).
\end{proposition}

\begin{corollary}[Freeing the null branch buys no predictions]
\label{cor:null}
Taking $c = 2a$, $a_i \equiv -a$ in Proposition~\ref{prop:gauge} shows that the set
$\{\,m + \phi_i + \phi_j : m \text{ free},\ \phi \text{ free}\,\}$ equals
$\{\,m_0 + \phi_i + \phi_j : m_0 \text{ fixed},\ \phi \text{ free}\,\}$. A free null branch and
a frozen one reach exactly the same composed predictions.
\end{corollary}

\begin{remark}[The gauge direction is the guidance-blind direction]
The move $u \mapsto u + \delta$, $\eps_1 \mapsto \eps_1 + \delta$ is the case $a_1 = 0$,
$a_2 = -\delta$, $c = \delta$ of Proposition~\ref{prop:gauge}. It leaves the loss exactly
invariant, and by \eqref{eq:Tij} it moves the sampled output by $-(w-1)\delta$. The
unidentified direction of the additive fit and the amplified direction of the sampler are the
same direction. Any variant is therefore incomplete without one of the three fixes already on
the table (frozen null in the guidance term, an anchor penalty
$\mu\norm{\epsth(\cnull) - \epsfr(\cnull)}^2$, or explicit per-concept deltas).
\end{remark}

\paragraph{Consequence for the comparison.}
The reachable set \eqref{eq:additive} at a fixed state depends on the prompts only through
which vector is called $\phi_i$. Renaming the branch prompt from \texttt{"a dog"} to
\texttt{"a dog, two animals"} to \texttt{"dog and"} renames a free vector. So, whenever the
per-concept functions are unconstrained, variants (a), (b) and (d-lora) have \emph{identical}
reachable sets, hence identical optima for \eqref{eq:dsm}. They differ in initialisation and in
conditioning number, not in capacity.

\section{The step-0 floor}
\label{sec:floor}

\paragraph{Why step 0 is exactly analysable.}
At the first sampling step the cached latent is bit-identical across every cell sharing a seed,
and so are the branch predictions: $\eps_\cnull$ is one tensor for all cells, and
$\eps(\text{``a baboon''})$ is the same tensor in every pair containing a baboon (verified, max
absolute difference $0.0$). So at that single state every per-concept adapter collapses to one
vector per concept, whatever its capacity, and \eqref{eq:additive} becomes a finite-dimensional
least-squares problem with no capacity assumption anywhere.

\paragraph{The floor.}
Let $E$ be the set of trained pairs, $y_{ij} \in \R^{d}$ the per-cell target at that state
($d = 4 \times 128 \times 128$ for SDXL at $1024^2$), and
\begin{equation}
  \boxed{\;
  R \;:=\; \min_{m,\ \phi_1,\dots,\phi_n}\ \sum_{(i,j)\in E}
  \norm{\,y_{ij} - m - \phi_i - \phi_j\,}^2 \;}
  \label{eq:floor}
\end{equation}
$R$ is a hard lower bound on the training loss of any variant whose free functions are
per-concept, at that state, with no assumption about the adapter's rank or architecture.

\begin{proposition}[$R$ is supported on the 2-core]
\label{prop:core}
If concept $i$ has degree $1$ in $G$, with unique edge $(i,j)$, then $\phi_i$ can be chosen to
fit that cell exactly and the edge contributes $0$ to \eqref{eq:floor}. Iterating, $R$ equals
the same minimum computed on the 2-core of $G$ (the maximal subgraph with minimum degree $2$).
\end{proposition}

\begin{remark}[Two cautions before the number is quoted]
With $1003$ cells over $294$ concepts (median degree $4$), the fit spends $294$ of $1003$
degrees of freedom, so an in-sample $R$ is deflated twice over: once by
Proposition~\ref{prop:core}, and once by the ordinary degrees-of-freedom count
($|E| - \mathrm{rank}$ residual dimensions per coordinate). The quantity to report is the
\emph{held-out} version: fit $(m,\phi)$ on a subset of pairs, evaluate on pairs excluded from
the fit. That is also the quantity the unseen-pair question actually asks.
\end{remark}

\paragraph{Per-variant floors.}
Write $\Delta_{ij} := y_{ij} - \big(\bar u + \bar f_i + \bar f_j\big)$ for the correction the
frozen composition still needs, and $\bar\Delta$ for its mean over cells.
\begin{align}
  R^{(a)} = R^{(b)} = R^{(d\ell)} &\;=\; R ,
  \label{eq:eqfloors}\\
  R^{(c)} &\;=\; \min_{\delta} \sum_{(i,j)\in E} \norm{\Delta_{ij} - \delta}^2
  \;=\; \sum_{(i,j)\in E} \norm{\Delta_{ij} - \bar\Delta}^2 ,
  \label{eq:cfloor}\\
  R^{(dt)} &\;=\; \min_{v \in \R^{d_\mathrm{text}}} \sum_{(i,j)\in E}
  \norm{\,y_{ij} - T^{(dt)}_{ij}(v)\,}^2 .
  \label{eq:dtfloor}
\end{align}
Line \eqref{eq:eqfloors} is Corollary~\ref{cor:null} plus the renaming argument;
\eqref{eq:cfloor} is a fit by one shared vector, minimised by the mean; \eqref{eq:dtfloor} has
no closed form. Both \eqref{eq:cfloor} and \eqref{eq:dtfloor} are minima over strict subfamilies of
\eqref{eq:additive}, so $R \le R^{(c)}$ and $R \le R^{(dt)}$. Neither $R^{(c)}$ nor $R^{(dt)}$
dominates the other: (c) is free in $\R^d$ but frozen in the concept branches, (d-tok) is free
only along the image of the text encoder but moves all three branches.

\paragraph{What the floor says about (c), against measurements already taken.}
By \eqref{eq:cfloor}, the fraction of correction energy (c) can remove is
$\norm{\bar\Delta}^2 / \E\norm{\Delta}^2$, the \emph{cross-pair} shared share at a common
state. The project's $19.6\%$ figure is not this quantity: it is cross-seed within one pair
(\texttt{scripts/rt\_shared\_component.py}, steps 0--2), and runs of the same pair should agree
more than runs of different pairs, so it is an upper bound at best. The matching measurement is
the D3b row, same noise and different prompt, which reads $+0.302$ over the first three steps
and $+0.006$ over steps 10--49; as an energy share that is of order $9\%$ early and nothing
late. The dose sweep then tested a shared correction directly (\texttt{mean\_others},
norm-matched, 128 cells) and scored AUC $0.059$ against the oracle's $0.387$ and a control band
reaching $0.047$. So $R^{(c)}$ is not merely worse than $R$; the variant it bounds has already
been close to falsified by measurement.

\paragraph{What would escape the floor.}
Nothing in (a)--(d). Escaping \eqref{eq:additive} requires either a read-out that sees both
concept embeddings at once,
\begin{equation}
  T_{ij} \;=\; u + f_i + f_j + h(e_i, e_j, \xs, \sigma),
  \qquad h \text{ symmetric, } h \not\equiv \text{additive},
  \label{eq:escape}
\end{equation}
which costs the property that the adapter never sees a joint prompt, or a sampler-side
correction (MCMC or Feynman-Kac), which is outside this family entirely.

\section{Is the plurality-conditioned composition the score of a density?}

\begin{proposition}
\label{prop:density}
$\;$Let $S(x) := \nabla \log p(x \mid c_1, \cpl) + \nabla \log p(x \mid c_2, \cpl) -
\nabla \log p(x \mid \cpl)$. Then $S = \nabla \log q$ with
\begin{equation}
  q(x) \;\propto\; \frac{p(x \mid c_1, \cpl)\; p(x \mid c_2, \cpl)}{p(x \mid \cpl)},
  \label{eq:q}
\end{equation}
provided $Z = \int p(x\mid c_1,\cpl)p(x\mid c_2,\cpl)/p(x\mid\cpl)\,dx < \infty$.
\end{proposition}
\begin{proof}
A sum and difference of gradients of logarithms is the gradient of the logarithm of the
corresponding product and quotient. Normalisability is the only extra condition.
\end{proof}

\paragraph{What the change of base measure actually assumes.}
By Bayes,
\begin{equation}
  p(x \mid c_1, c_2, \cpl)
  \;=\; \frac{p(x\mid c_1,\cpl)\,p(x\mid c_2,\cpl)}{p(x\mid \cpl)}
  \cdot \underbrace{\frac{p(c_1,c_2 \mid x,\cpl)}{p(c_1\mid x,\cpl)\,p(c_2 \mid x,\cpl)}
  \cdot \frac{p(c_1\mid\cpl)p(c_2\mid\cpl)}{p(c_1,c_2\mid\cpl)}}_{\text{the interaction term}},
  \label{eq:bayes}
\end{equation}
so $q$ of \eqref{eq:q} is the exact posterior when $c_1 \perp c_2 \mid (x, \cpl)$. That is
strictly weaker than the unconditional independence the current composition assumes, and weaker
in the direction that matters: it only has to hold among images that already contain two
separate objects, which is where the chimera mode has been conditioned away. The correction
term in \eqref{eq:bayes} is the plurality-conditioned version of the pointwise mutual
information whose gradient is $r_t$ in this project's own notation.

\paragraph{Two limits on that argument.}
First, Proposition~\ref{prop:density} is a clean-data statement. At $\sigma > 0$ the score of
the noised $q$ is not $S$ diffused, so an exact (b) composition is the score of $q$ only at
$\sigma = 0$. Objective \eqref{eq:dsm} sidesteps this by regressing per level
(Proposition~\ref{prop:minimiser}), for (a) as much as for (b). Second, the DSM optimum is the
projection of $\nabla \log p_{\mathrm{co},\sigma}$ onto the family, not $q$. So (b)'s density
argument justifies its inductive bias and its starting point, not its optimum, and by
\eqref{eq:eqfloors} its optimum coincides with (a)'s.

\section{Is \texorpdfstring{$c \oplus \cpl$}{prompt concatenation} well defined on SDXL?}
\label{sec:concat}

\texttt{encode\_prompt\_sdxl} (\texttt{poe\_repair/\_sdxl/runtime.py:176--211}) tokenises the
whole prompt to a fixed $77$ positions in both tokenizers with \texttt{padding="max\_length"},
takes \texttt{hidden\_states[-2]} from each encoder, and concatenates them on the channel axis:
\begin{equation}
  \mathrm{enc}(p) = \Big(\underbrace{[\,h^{(1)}_{-2}(p) \,\Vert\, h^{(2)}_{-2}(p)\,]}
  _{\in\,\R^{77\times 2048}},\ \underbrace{\mathrm{pool}^{(2)}(p)}_{\in\,\R^{1280}}\Big).
\end{equation}
Consequently:
\begin{itemize}
  \item \textbf{Embedding addition is not the operation.}
    $\mathrm{enc}(p_1) + \mathrm{enc}(p_2)$ aligns token slot $k$ of one prompt onto slot $k$
    of the other and sums sixty-odd padding and EOS embeddings, producing a sequence no string
    maps to. The pooled vector is worse: it is a single summary of one string, fed to the
    time-embedding path through \texttt{added\_cond\_kwargs["text\_embeds"]}, and a sum of two
    summaries summarises nothing.
  \item \textbf{String concatenation before tokenisation is the operation}, and it is well
    defined: $c \oplus \cpl$ means one string (\texttt{"a dog, two animals"}) encoded once. The
    $77$-token budget is nowhere near binding for these prompts.
\end{itemize}
The same reading applies to (d): \texttt{"dog and"} is a string, and the LoRA in \eqref{eq:vdl}
is not localised to its last token.

\paragraph{Cache consequence.}
The cached branch tensors \texttt{eps\_a\_raw}, \texttt{eps\_b\_raw}, \texttt{eps\_uncond} are
keyed on the current prompts. Variants (b) and (d) change those prompts and invalidate them;
(a) and (c) do not.

\section{Feasibility}

\paragraph{Cost per training step is lower than the objective it replaces.}
Both \eqref{eq:distill} and \eqref{eq:dsm} need three adapted UNet forwards per sample. The
distillation target additionally needs the frozen joint-prompt prediction (today read from
cache, which is why the cache exists). The DSM target is $z$, which is free. So per step,
\eqref{eq:dsm} is strictly cheaper and removes a cache dependency.

\paragraph{Target noise is smallest exactly where the correction lives.}
Since $z = (\xs - x)/\sigma$,
\begin{equation}
  \mathrm{Var}[\,z \mid \xs\,] \;=\; \frac{1}{\sigma^2}\,\mathrm{Var}[\,x \mid \xs\,]
  \;\xrightarrow[\ \sigma \to \infty\ ]{}\; \frac{\mathrm{Var}_{p_{\mathrm{co}}}[x]}{\sigma^2}
  \;\to\; 0 ,
  \label{eq:targetnoise}
\end{equation}
while at $\sigma \to 0$ it stays $O(1)$ per coordinate. The regression target is therefore
\emph{least} noisy at large $\sigma$, which is the first few sampling steps, which is the
window this project has measured the correction to occupy. Gradient variance is the usual
objection to DSM against a distillation target, and \eqref{eq:targetnoise} says the objection
is weakest in the only window that matters. \emph{Confidence: the limits are exact; the
magnitude at the $\sigma$ values of the DDIM grid used here is unverified.}

\paragraph{The binding constraint is the data, not the loss.}
$p_{\mathrm{co}}$ does not exist in this repository. It requires real images with both concepts
present as separate objects, over a $294$-concept vocabulary whose pairs were selected to be
confusable species. Every cached $\xs$ is a PoE sampling-trajectory latent rather than a noised
real image, so no existing cache is the substrate for \eqref{eq:dsm} under any variant.

\paragraph{Does the floor transfer to the new objective?}
Not as a bit-exact identity. Step-0 exactness relies on one shared cached latent; under
\eqref{eq:dsm}, $z$ is drawn per sample, so states differ across cells even at large $\sigma$.
What survives is the statistical version: as $\sigma$ grows the noised distributions overlap,
so one $\phi_i$ must serve every pair containing concept $i$, and \eqref{eq:additive} binds in
the limit. Separately, the irreducible term $\E\norm{z - \E[z\mid\xs]}^2$ is
\emph{family-independent}, so it cancels from every difference of floors. Comparisons between
variants under \eqref{eq:dsm} are therefore meaningful even though the absolute loss is not.

\paragraph{Summary table.}
\begin{center}
\begin{tabular}{llll}
\toprule
Variant & Free objects at one state & Floor & Reuses current prompt cache \\
\midrule
(a) plain        & $\{\phi_i\}_{i=1}^{n}$, $m$        & $R$                 & yes \\
(b) plurality    & $\{\phi_i\}_{i=1}^{n}$, $m$        & $R$ (equal to (a))  & no \\
(c) null-only    & one $\delta \in \R^d$              & $R^{(c)} \ge R$     & yes \\
(d-lora)         & $\{\phi_i\}_{i=1}^{n}$, $m$        & $R$ (equal to (a))  & no \\
(d-tok)          & one $v \in \R^{d_\mathrm{text}}$   & $R^{(dt)} \ge R$    & no \\
\bottomrule
\end{tabular}
\end{center}

\section{Recommendation, and the number that would overturn it}

Train (a) with the null branch anchored,
\begin{equation}
  \mathcal{L} \;=\; \mathcal{L}_{\mathrm{DSM}}
  \;+\; \mu \, \E\,\norm{\,\epsth(\xs,\sigma,\cnull) - \epsfr(\xs,\sigma,\cnull)\,}^2 ,
  \label{eq:recommended}
\end{equation}
the anchor fixing the gauge of Proposition~\ref{prop:gauge} at no extra forward pass. (a)
changes exactly one axis from the current setup (the target), so its comparison against the
existing distillation run is clean; every other variant changes the target and the conditioning
together. By \eqref{eq:eqfloors} the other three cannot beat it on capacity: two tie, two are
strictly smaller.

\paragraph{The overturning measurement.}
Compute the held-out-pair version of \eqref{eq:floor} at step 0 against the existing cached
targets, as a fraction of the correction energy,
\begin{equation}
  \rho \;:=\; \frac{R_{\text{held-out}}}{\sum_{(i,j)} \norm{\Delta_{ij}}^2},
  \qquad
  \text{support if } \rho < 0.5,\quad \text{kill if } \rho \ge 0.5,
  \label{eq:bar}
\end{equation}
with the bar written into the script before the number is produced. If $\rho$ is large, no
per-concept variant can reach the target however it is trained, and the work belongs on the
sampler side or on the interaction read-out \eqref{eq:escape}. This uses only the existing
cache and the existing target, so it costs no GPU time beyond reading tensors.

\end{document}
